% ══════════════════════════════════════════════════════════════════════════════
% TEMPLATE: Gruppenübung
%
% Altersgruppe:  welpen | pubertaet | erwachsen
% Schwierigkeit: 1 (leicht) bis 5 (sehr schwer)
% Bild (optional): images/gruppenuebungen/ALTERSGRUPPE/dateiname.png
%
% Hinweis: Gruppenübungen haben zwei Bilder –
%   Bild 1 (Seitenansicht): zeigt Trainerin und Hund
%   Bild 2 (Draufsicht):    zeigt den Übungsablauf als Diagramm
% ══════════════════════════════════════════════════════════════════════════════

\begin{uebung}{TITEL}{ALTERSGRUPPE}{KATEGORIE}{SCHWIERIGKEIT}

  % Bild Seitenansicht (Zeile entfernen wenn kein Bild vorhanden)
  \uebungsbild{images/gruppenuebungen/ALTERSGRUPPE/DATEINAME_seite.png}

  \uebungsbeschreibung{%
    KURZE EINLEITUNG was in dieser Gruppenübung trainiert wird
    und warum die Gruppe dabei wichtig ist.\\[6pt]
    \textbf{Aufbau:}
    \begin{itemize}[leftmargin=*, itemsep=2pt]
      \item Teilnehmer: ANZAHL PERSONEN, ANZAHL HUNDE
      \item Aufstellung: BESCHREIBUNG DER STARTPOSITION
      \item Material: BENÖTIGTES MATERIAL (z.B. Leine, Clicker, Leckerlis)
    \end{itemize}

    \vspace{6pt}
    \textbf{Ablauf:}
    \begin{enumerate}[leftmargin=*, itemsep=2pt]
      \item SCHRITT 1
      \item SCHRITT 2
      \item SCHRITT 3
      \item SCHRITT 4
    \end{enumerate}
  }

  \uebungsvariationen{%
    \begin{itemize}[leftmargin=*, itemsep=4pt]
      \item \textbf{VARIATION 1:} BESCHREIBUNG
      \item \textbf{VARIATION 2:} BESCHREIBUNG
      \item \textbf{VARIATION 3:} BESCHREIBUNG
    \end{itemize}
  }

  \uebungstrainer{%
    \begin{itemize}[leftmargin=*, itemsep=4pt]
      \item \textbf{Gruppenmanagement:} HINWEIS ZUR KOORDINATION DER GRUPPE
      \item \textbf{Sicherheit:} HINWEIS ZU SICHERHEITSASPEKTEN
      \item WEITERER TRAINER-HINWEIS
      \item WEITERER TRAINER-HINWEIS
    \end{itemize}
  }

\end{uebung}
